\documentclass[11pt]{article}
\usepackage[a4paper,margin=25mm]{geometry}
\usepackage{amsmath,amssymb,amsthm,microtype,hyperref}
\newtheorem{theorem}{Theorem}\newtheorem{lemma}{Lemma}
\newcommand{\E}{\mathbb E}\newcommand{\Pp}{\mathbb P}
\title{A renewal identity for the one-sided conditional-noise case}
\author{}\date{28 October 2026}
\begin{document}\maketitle

The count proof estimates several crossing probabilities separately. When
the lower conditional law has no noise, those crossings are successive
renewals of one positive random walk. Their sum has a simple sign, with no
Poisson comparison and no parameter certificate. This does not yet control
the finite binomial weights or cover two noisy conditional laws.

\begin{theorem}\label{thm:renewal}
Let $D_1,D_2,\ldots$ be iid, with $\E D_1=1$ and
\[
       0<m\le D_1\le M,\qquad b:=1-(M-m)\ge0.
\]
Set $T_0=0$, $T_j=D_1+\cdots+D_j$,
$N(t)=\max\{j:T_j\le t\}$ for $t\ge0$, and let
$A(t)=t-T_{N(t)}$ be the age of the last renewal. Then
\[
 g(t):=t-\E N(t)\ge \E\min\{A(t),b\}\ge0.                 \tag{1}
\]
Both endpoints in the support assumption are permitted. Atoms at the
evaluation point count as renewals; the next crossing is strictly above it.
\end{theorem}
\begin{proof}
Write $\overline F(x)=\Pp(D_1>x)$ and
\[
 h(x)=\overline F(x)-\E(D_1-x)_+,\qquad x\ge0.
\]
If $0\le x<m$, then $h(x)=x$ and $\overline F(x)=1$.
If $x\ge m$, then
$(D_1-x)_+\le(M-m)\mathbf1_{\{D_1>x\}}$, so
\[
 h(x)\ge b\overline F(x)\ge\min\{x,b\}\overline F(x).
                                                               \tag{2}
\]
The latter lower bound also holds for $x<m$.

The renewal equation for $q(t)=\E N(t)$ is
$q(t)=\Pp(D_1\le t)+\E[q(t-D_1)\mathbf1_{\{D_1\le t\}}]$.
Substitute $q(t)=t-g(t)$, use $\E D_1=1$, and obtain
\[
 g(t)=h(t)+\E[g(t-D_1)\mathbf1_{\{D_1\le t\}}].             \tag{3}
\]
Every increment is at least $m>0$, so iteration is a finite operation:
\[
 g(t)=\sum_{j=0}^{\lfloor t/m\rfloor}
       \E[h(t-T_j)\mathbf1_{\{T_j\le t\}}].                 \tag{4}
\]
For each $j$, independence of the next increment identifies
$\overline F(t-T_j)$ with the conditional probability of
$T_j\le t<T_{j+1}$. These disjoint events partition the sample space.
Combining (2) and (4) gives (1). This proves the claim without an optional
stopping theorem, a limiting renewal theorem, or a non-lattice assumption.
\end{proof}

\section*{The exact iid application}
Suppose $Y=B+U$, $\Pp(B=1)=p$, $\E(U\mid B)=0$, and $U_0=0$.
Let $U_1\in[-a,A]$, where $a<1$ and $a+A\le1$.
Use independent copies of $D=1+U_1$, so $m=1-a$, $M=1+A$ and
$b=1-a-A$. If
\[
                         0\le t<n(1-a),                    \tag{5}
\]
then $N(t)\le n-1$. Conditional on $j$ successes, the raw sum of $n$
copies of $Y$ has exactly the law of $T_j$. Therefore, with $F_B$ the
$\mathrm{Bin}(n,p)$ CDF,
\[
             \Pp(Y_1+\cdots+Y_n\le t)=\E F_B(N(t)).        \tag{6}
\]
Indeed the indicators $\mathbf1_{\{T_j\le t\}}$ equal
$\mathbf1_{\{j\le N(t)\}}$ on this one common probability space.
This is an identity for an actual iid marginal, not a relaxation.

The acquired right strip $a\le.072$, $A=.6$ satisfies $b\ge.328$;
the left strip $a\le.068$, $A=.7$ satisfies $b\ge.232$.
For $n\ge42$, $p\in[.25,.53]$, $\lambda=n\E U^2\le.616$ and
$|H|\le1.84$, the raw threshold
$t=np+H\sqrt{p(1-p)+\lambda/n}$ satisfies (5). For example
$\sqrt{p(1-p)+\lambda/n}<.515$ gives
$0<np-.948\le t\le.53n+.948<.928n$.
Thus (6) and $\E N(t)\le t$ hold throughout this entire noise range
in the stated one-sided case. No claim that $U_0=0$ follows from the
acquisition hypotheses is made.

\section*{What the identity does and does not remove}
Put $v=p(1-p)$, $s=\E U^2$, $\rho=\E|Y-p|^3/(v+s)^{3/2}$,
$\rho_B=(1-2v)/\sqrt v$, and $t=k+u$. Write $\pi_k=\Pp(B_n=k)$,
$m_k=\sqrt n\,\pi_k$, and define the exact discrete remainder
\[
 R_k(j)=F_B(j)-F_B(k)-\pi_k(j-k).
\]
For the direct comparison used in the preceding research notes, (6) gives
\begin{align*}
 \Delta_C={}&-m_k g(t)+\sqrt n\,\E R_k(N(t))+m_k u\\
 &-\sqrt n\left[
 \Phi\!\left(\frac{t-np}{\sqrt{n(v+s)}}\right)
 -\Phi\!\left(\frac{k-np}{\sqrt{nv}}\right)\right]
 +C(\rho_B-\rho).                                           \tag{7}
\end{align*}
There is no sum of separately bounded incoming probabilities in (7).
The leading unweighted contribution is nonpositive for every phase and
every noise size satisfying the support assumptions. But (1) alone does
\emph{not} dominate the binomial remainder, the normal correction, or the
cubic correction in (7). A bound relative to arbitrarily small noise must
retain their joint dependence. Consequently this note removes no dependency
from the already accepted general comparison theorem. It provides a
rigorous alternative formula and a useful test case for a future replacement.

The width condition in Theorem~\ref{thm:renewal} cannot simply be dropped.
Take $D=9/10$ with probability $19/21$ and $D=39/20$ with probability
$2/21$. Its mean is one, but at $t=9/10$,
$\E N(t)=19/21>9/10$, with excess $1/210$.
The support width is $21/20$. Also, replacing $\E N(t)$ by $t$ would
be wrong even for the deterministic increment $D=1$, when
$g(t)=t-\lfloor t\rfloor$.

\section*{Evidence and scope}
The proof above is elementary and covers all distributions with the stated
support and mean. The companion \texttt{renewal-check.py} uses exact finite
convolutions to check (3)--(4), the age bound, atom collisions and the two
false extensions. These tests are controls on the derivation, not a substitute
for its distribution-free argument. No full Berry--Esseen comparison,
finite-base bound, or global theorem is asserted here.
\end{document}
